\section{The entire triangular source-curvature append step above first weight eight}
This is an unbounded CONDITIONAL propagation theorem. It does not prove
the triangular target at its low-first-weight entrance layer, and neither
the triangular target nor the original conjecture is declared complete.

For an actual normalized nonrectangular sorted arm let $\kappa,\delta$
be its first missing weight and leading source multiplicity. Put
\[
 E=1-z+(2z-1)F,\quad D_i=(F_i-F_{i-1})^+\ (i\ge2),
 \quad D_0=D_1=0,\quad \ell=\min\operatorname{argmax}D,
 \quad B(d)=\max\{1,\tbinom{d+1}{2}-1\}.
\]
Write $\mathsf{TC}(F)$ for the assertion
\[
          \Delta E_n\ge B(\delta)
                     \qquad(\kappa+3\le n\le\ell).       \tag{TC}
\]
The source is zero below $\kappa$, $1\le\delta\le\kappa$, and
its increments through $\ell$ are nonnegative by actual single-arm
positive-derivative doubling. All sequences are zero extended.

\begin{theorem}[Complete large-first-weight append propagation]
Suppose $\kappa\ge8$ and $\mathsf{TC}(F)$ holds for this OLD arm.
For EVERY legal append $t$, the NEW arm satisfies $\mathsf{TC}$.
No bound on the number of parts, their sizes or the initial run is imposed.
\end{theorem}
\begin{proof}
Use hats for the new arm, fix $\widehat\kappa+3\le n\le\widehat\ell$,
and use the exact recurrence
\[
 \Delta\widehat E_n
   =\sum_{i=n-t}^{n-1}\Delta E_i
          +\mathbf1_{n=t+1}-\mathbf1_{n=t+2}.           \tag{1}
\]
If $n-1>\ell$, the peak-crossing source payment in
\texttt{agent\_postcount\_source\_mode\_payment.tex} gives
$\Delta\widehat E_n\ge\max D-1\ge B(\widehat\delta)$ for
$\kappa\ge8$. Its proof contains the old derivative peak inside
the window and pays the possible negative unit by the exponential
first-missing mass. That branch is unconditional and uses no OLD TC.
It remains to treat $n-1\le\ell$. All summands in (1) are then
nonnegative. There are exactly three legal first-missing transitions.

\paragraph{Stationary transition $t=\kappa-1$.}
Here $\widehat\kappa=\kappa$ and $\widehat\delta=1$.
The negative unit would be at $n=\kappa+1$, outside the required range.
At $n=\kappa+3$, the sum in (1) is exactly
$E_{\kappa+2}-E_3=E_{\kappa+2}$. The early unscaled dyadic profile
gives $E_{\kappa+2}\ge2$, paying $B(1)=1$.
For $n\ge\kappa+4$, the single summand at $n-1$ is at least
$B(\delta)\ge1$ by OLD TC. The other summands and possible unit
are nonnegative, proving this transition.

\paragraph{Preserving growth $t>\kappa$.}
Here $\widehat\kappa=\kappa+1$, $\widehat\delta=\delta$,
and $n-1\ge\kappa+3$. The summand at $n-1$ supplies OLD TC.
Only the negative unit at $n=t+2$ could reduce the payment.
At that index the required range forces $t\ge\kappa+2$,
and the window $[2,t+1]$ also contains the distinct birth rank
$\kappa$, whose increment is exactly $\delta\ge1$.
This additional term pays the negative unit. Thus the growth
transition preserves (TC), including its unit dent.

\paragraph{Tied growth $t=\kappa$.}
Here $\widehat\kappa=\kappa+1$ and $\widehat\delta=\delta+1$.
Both isolated units are before the required range. At its first rank
$n=\kappa+4$, (1) is
\[
                  E_{\kappa+3}-E_3=E_{\kappa+3}.        \tag{2}
\]
We claim the following unconditional coefficient bound:
\[
             E_{\kappa+3}\ge B(\delta+1)
                       =\frac{\delta(\delta+3)}2
                              \qquad(\kappa\ge8).       \tag{3}
\]
For $1\le\delta\le5$, the multiplicity-scaled early dyadic bound
gives $E_{\kappa+3}\ge4\delta$, because $3\le\kappa-\delta$;
this dominates the right side of (3). For $\delta\ge6$, use
$E_{\kappa+3}\ge[z^3]M_\delta$. The exact profile has
$[z^3]M_6=44\ge27$. For every $d\ge6$ its recursion gives
\[
 [z^3]M_{d+1}-[z^3]M_d
    =[z^2]M_d+[z]M_d+d\ge d+2,
\]
since $[z]M_d=\binom d2-1\ge2$. These are precisely at least
the increments of $d(d+3)/2$, proving (3) by induction.

For $n\ge\kappa+5$ and $\delta\ne2$, the two distinct window
summands at $n-1,n-2$ both fall in OLD TC's range. They pay because
\[
                  2B(d)\ge B(d+1)\qquad(d=1\text{ or }d\ge3).
                                                               \tag{4}
\]
The exceptional multiplicity $\delta=2$ needs three old payments:
$3B(2)=6\ge B(3)=5$. These are present for $n\ge\kappa+6$,
at $n-1,n-2,n-3$. At the one intervening rank $n=\kappa+5$,
the exact sum is $E_{\kappa+4}-E_4=E_{\kappa+4}$.
The scaled early dyadic bound gives $E_{\kappa+4}\ge16$, since
$4\le\kappa-2$, and therefore pays $B(3)=5$.
This covers every tied-growth rank and completes the proof.
\end{proof}

\begin{corollary}[Only the low entrance layer is needed for this target]
If (TC) is proved for EVERY actual arm with first missing weight eight,
then it holds for every actual arm with first missing weight at least
eight. Equivalently a least-part counterexample to (TC) in that region
must enter weight eight from an old arm of first missing weight seven.
\end{corollary}
\begin{proof}
Induct on the number of non-one parts, using that a legal append
preserves the first missing weight or increases it by exactly one.
An arm above eight has predecessor at least eight. Stationary
extensions at eight are covered by the same theorem. Its first entrance
from seven is included in the assumed base family.
\end{proof}

This is a reduction to a BOUNDED FIRST-WEIGHT family, not an exhaustive
finite check: initial runs, capacities, counts and comparison ranks in
the weight-eight entrance family are still unbounded. The subsequent
\texttt{triangular\_source\_curvature\_complete.tex} closes this entrance
family and ALL smaller first weights by a separate rank reduction and
two complete independently rebuilt finite jet certificates.
The full triangular theorem alone does not pay the entire growing
two-source convolution at later overlaps.

The exact actual diagnostic
\texttt{probe\_triangular\_source\_induction.py/.json} checks 50000
append contexts and 2166342 target comparisons, including 1461270
large-first-weight pre-old-peak comparisons and all three local
boundary identities. No failure was found. This is diagnostic evidence
only; the conditional unbounded proof is the argument above.
